\chapter{Introduction}
\pagenumbering{arabic}\hspace{3mm}

Before the Transformer architecture \cite{attention_paper} was introduced, neural networks had to read text one word at a time. This step-by-step approach was slow and made it difficult for models to remember the beginning of a long paragraph by the time they reached the end. The Transformer solved this by processing the entire sequence at once using a mechanism called self-attention. Inside the network, it transforms each word into mathematical representations known as queries, keys, and values. This allows the model to capture deep context and relationships, no matter how far apart the words are. Over time, this incredibly flexible foundation was adapted into three main types of models to handle different kinds of problems. Encoder-only models, like BERT \cite{bert_paper}, are built to read text in both directions to get a comprehensive understanding of the context, making them excellent for sequence classification and analysis tasks. Encoder-decoder models, such as T5 \cite{t5_paper}, use one part to process the input and another to generate a completely new sequence, which is ideal for tasks like translation. Finally, decoder-only models \cite{gpt3_paper} focus strictly on predicting the next word based on the text that came before it. Because they scale so seamlessly to handle massive amounts of data, these decoder-only architectures are the primary engines driving today's generative AI systems.

When these decoder-only Transformers are scaled to billions of parameters, they become Large Language Models (LLMs) \cite{gpt3_paper}. At this scale, models such as OPT \cite{opt_paper}, LLaMA \cite{llama_paper}, and Mistral \cite{mistral_paper} begin to exhibit emergent capabilities beyond their explicit training objectives \cite{gpt3_paper}. However, the models' capacity has grown exponentially, outpacing the hardware technology intended to support it \cite{zeroquant_paper, smoothquant_paper, awq_paper}. Over the last five years, parameter counts have increased from hundreds of millions to hundreds of billions, while GPU memory capacities have barely doubled \cite{llm_int8_paper, smoothquant_paper}. A GPT-3 scale model in FP16 needs around 350 GB just to hold its weights \cite{gptq_paper, zeroquant_paper}. Serving it requires multiple high-end accelerators per replica, before any room is left for activations, KV-cache, or batching \cite{llm_int8_paper}. The same gap shows up at smaller scales too: a 7B model in FP16 already exceeds the memory of most consumer GPUs and almost all mobile NPUs \cite{awq_paper, gptq_paper}.

This makes inference, not training, the practical bottleneck for deploying these models \cite{smoothquant_paper, outlier_suppression_paper}. Training is a one-time cost, whereas inference cost is paid on every user input, strictly determining whether a model can physically run on a given device. Thus, reducing the compute and memory cost of inference is essential to making these models usable outside of data centers---on laptops, phones, and embedded hardware where the memory budget is limited \cite{llm_int8_paper, gptq_paper, awq_paper}.

\section{Background}
\textbf{Quantization} which involves replacing high-precision floating-point tensors with low-bit integers is one of the few techniques that addresses memory footprint, memory bandwidth, and arithmetic throughput at the same time \cite{gptq_paper, zeroquant_paper}. It also maps to hardware that is already deployed at scale: NVIDIA Tensor Cores, Intel AMX,and Qualcomm DSPs all provide native INT8 GEMM kernels with roughly twice the throughput and half the memory traffic of FP16 \cite{llm_int8_paper, smoothquant_paper}.There are two broad ways to quantize a model. \textbf{Quantization-aware training (QAT)} retrains the model with simulated low precision and recovers accuracy well, but its cost is comparable to pre-training itself. At LLM scale, this is not a practical option because performing backpropagation requires storing the full computation graph, gradients, and optimizer states in memory, demanding massive compute clusters. \textbf{Post-training quantization (PTQ)} keeps the trained weights frozen and quantizes them in a single calibration pass on a small held-out sample \cite{gptq_paper, awq_paper, zeroquant_paper}. PTQ hence is the setting of this thesis.

The standard form of integer quantization used throughout this work is symmetric uniform quantization \cite{smoothquant_paper, llm_int8_paper}:
\begin{equation}
\bar{X}^{\text{INT8}} = \left\lceil \frac{X^{\text{FP16}}}{\Delta} \right\rfloor, \qquad \Delta = \frac{\max(|X|)}{2^{N-1} - 1}
\end{equation}

with $N = 8$ bits. The granularity at which $\Delta$ is computed defines the scheme. \textit{Per-tensor} uses one step size for the whole matrix. \textit{Per-token} uses one per row of the activation. \textit{Per-channel} uses one per output channel of the weight. Coarser schemes are cheaper and map directly to INT8 GEMM kernels. Finer schemes give each row or channel its own range but add scaling work that has to be fused into the matmul \cite{smoothquant_paper, zeroquant_paper}.

A naive expectation when moving weights from FP16 to INT8 is a 50\% drop in peak GPU memory at inference. In practice this does not hold. This has been verified in the early stages of this thesis through various experiments measuring the memory saving as we reduce precision from FP32 to FP16.  As the model grows, the peak memory of a forward pass is dominated not by the weights but by the \textbf{activations}, the intermediate tensors computed inside attention and feed-forward layers, including the Q/K/V projections, the attention probability matrix, the FFN hidden state, and the KV-cache. Quantizing only the weights gives a memory saving that gets smaller as the model gets larger, because the un-quantized activations grow to dominate the budget. This is the \textbf{W8A8} referenced in LLM.int8() \cite{llm_int8_paper}, SmoothQuant \cite{smoothquant_paper}, and this work, representing both weights and activations in INT8. It is also what is required by the INT8 GEMM kernels for both operands to be INT8, in order to invoke the tensor cores of the GPU \cite{smoothquant_paper, llm_int8_paper}. 

However, applying W8A8 PTQ directly to massive LLMs typically results in catastrophic accuracy degradation due to the emergence of extreme activation outliers. Two properties of these outliers are central to the rest of this thesis \cite{llm_int8_paper}:
\begin{description}
    \item[Outliers are sparse:] Roughly 0.1\% of feature dimensions, but those channels carry information that the model needs.
    \item[Outliers persistence within a channel:] If a channel is an outlier channel, it is large for \textbf{every} token.
\end{description}.
In per-tensor activation quantization, a single global scale factor $\Delta$ is dictated by the largest outlier. This severely penalizes typical channels with magnitude $m_i$; when the global maximum is $m$, the available integer levels shrink to $2^8 \cdot \frac{m_i}{m}$, sometimes leaving only two or three usable bins \cite{smoothquant_paper}. This dynamic causes catastrophic accuracy collapse: on OPT-175B, naive W8A8 degrades zero-shot average accuracy from 71.6\% to 32.3\%, nearing the random-guessing floor \cite{smoothquant_paper}.

Two prior approaches attempted to resolve this. \textbf{LLM.int8()} \cite{llm_int8_paper} isolates outlier channels in FP16 and routes the remaining matrix multiplications through INT8. While this preserves accuracy, the resulting mixed-precision kernel breaks the monolithic INT8 execution path, making it in practice slower than the FP16 baseline it aimed to accelerate \cite{smoothquant_paper}. Alternatively, \textbf{ZeroQuant} \cite{zeroquant_paper} employs highly fine-grained quantization---using per-token activations and group-wise weights via custom CUDA kernels. While effective for models up to roughly 20B parameters, it fails to sustain accuracy at the 175B scale.

SmoothQuant \cite{smoothquant_paper} is the foundational method this thesis builds upon, offering a fundamentally different solution to the outlier problem. Building on the observation that outliers are stable per channel but highly variable across channels \cite{llm_int8_paper}, the authors recognized that per-channel activation quantization is mathematically ideal, but practically impossible to fuse inside standard INT8 matrix multiplications.

To resolve this, SmoothQuant introduces a mathematically equivalent transformation offline: it mathematically "smooths" the activations by migrating the quantization difficulty off the activations and into the weights. Because the original weight distributions are relatively flat, they can absorb this extra variance without significant degradation. By smoothing the activations before inference, SmoothQuant allows the model to utilize highly efficient, hardware-friendly INT8 operations. Across massive architectures like OPT, BLOOM, and GLM-130B, this method matches FP16 zero-shot accuracy to within a fraction of a percent while delivering substantial speedups and memory savings \cite{smoothquant_paper}. Because it remains the strongest training-free, fully-INT8 baseline available, it serves as the natural starting point for the optimization strategies proposed in this work.

\section{Motivation}

SmoothQuant provides a clean, training-free mechanism to handle activation outliers by migrating their quantization difficulty into the weights. However, its baseline implementation relies on two coarse design choices: a single per-channel calibration statistic (the absolute maximum) and a single migration strength ($\alpha$) shared uniformly across the entire network.A careful analysis of the actual outlier distributions in Large Language Models reveals that these distributions are significantly more structured than a global static configuration can capture. This thesis is motivated by three distinct empirical gaps in the baseline recipe, leading to three principled refinements.

\subsection{Per-Channel Weight Quantization}

The fundamental premise of SmoothQuant is to shift the variance of outlier channels away from the activations and into the weights \cite{smoothquant_paper}. Consequently, after smoothing, the adjusted weight tensor carries significantly more variance than the original weights.

If these adjusted weights are quantized using a single, per-tensor scale (as in the baseline configurations), channels with typical magnitudes are inevitably squeezed into a very small number of effective integer levels. The natural countermeasure to this induced variance is to apply per-channel quantization to the weights \cite{llm_int8_paper}. Because the output-channel dimension is an outer dimension of the matrix contraction, per-channel weight scaling does not interfere with the integer arithmetic of standard INT8 GEMM kernels \cite{smoothquant_paper}.

\subsection{Quantile-Based Smoothing}

The second gap lies in the lower level calibration step. The baseline formula calculates the scale using the absolute maximum of each activation channel over a calibration buffer. While mathematically simple, this makes the entire quantization process hypersensitive to single-token anomalies.

In practice, within a typical non-outlier channel, the vast majority of tokens sit near the median, but a few isolated tokens can spike well above it. Because the baseline method uses the absolute maximum, a single anomalous spike in the calibration data dictates the scaling factor for that channel across every subsequent forward pass during deployment. The SmoothQuant authors themselves run into this on GLM-130B and patch it by clipping the top 2\% of tokens before computing static quantization step sizes \cite{smoothquant_paper}. The clip is effective but ad hoc since it lives at the quantization step rather than the smoothing step. A more robust approach is to replace the absolute maximum with a high empirical quantile (e.g., the 99th or 99.9th percentile). This ignores isolated, single-token noise while accurately capturing the true scale of the channel's distribution, preventing calibration noise from degrading inference accuracy.

\subsection{Per-Layer Migration Strength}

Finally, the severity of activation outliers is not uniform across the depth of the network. As the outlier severity varies across different layers. Across OPT-2.7B, OPT-6.7B, and OPT-13B, the ratio between the most outlier-heavy layer and the cleanest layer of the same model grows. A layer with mild activation outliers requires a small ($\alpha$) that leaves the weights largely untouched, while a layer with very heavy outliers benefits from a larger ($\alpha$) that shifts more of the difficulty onto its weights. 

The baseline method applies a single global migration strength ($\alpha$) to the entire network (0.5 for OPT and BLOOM, 0.75 for GLM-130B) \cite{smoothquant_paper}. Because the outlier severity profile is static and can be measured offline, a natural refinement is to replace the global $\alpha$ with a dynamic, per-layer $\alpha^{(l)}$. By assigning a smaller $\alpha$ to layers with mild outliers and a larger $\alpha$ to layers with heavy outliers. Since this is an offline, one-shot modification of the calibration step, it does not add any extra cost during the inference.

\medskip
These three refinements share a common observation. The baseline SmoothQuant configuration treats the network as if it were statistically uniform across weights, channels, and layers. The actual outlier distribution observed inside LLMs is not. Per-channel weight quantization, quantile-based smoothing, and per-layer migration strength each address one specific mismatch in this assumption. All three live entirely at the offline calibration step and therefore add no extra cost at deployment. The next section turns these three observations into the concrete objectives this thesis aims to verify.

\section{Objectives}

The primary aim of this thesis is to improve the post-training quantization of Large Language Models by refining both the offline calibration stage and the structural smoothing stage. This work pursues the following four targeted objectives:

\subsection{Integrating Quantile Calibration with Per-Channel Weights}
The first objective is to pair per-channel weight quantization with a robust, quantile-based smoothing statistic, completely replacing the noise-sensitive absolute maximum used in prior work. This integrated configuration will be evaluated on mid-sized architectures (OPT-125M, OPT-1.3B, and OPT-2.7B)\cite{opt_paper}. Its performance will be benchmarked against the standard max-based SmoothQuant baselines using Wikitext-2 Perplexity (PPL) and zero-shot accuracy across seven standard downstream tasks: LAMBADA \cite{lambada_dataset}, HellaSwag \cite{hellaswag_dataset}, PIQA \cite{piqa_dataset}, WinoGrande \cite{winogrande_dataset}, OpenBookQA \cite{openbookqa_dataset}, RTE \cite{rte_dataset}, and COPA \cite{copa_dataset}.

\subsection{Developing a Per-Layer Migration Strength}
The second objective addresses the inefficiency of a rigid, globally shared migration strength ($\alpha$). This involves deriving a per-layer heuristic, $\alpha^{(l)}$, that takes into account the measured outlier severity of each layer during the calibration step. This smoothing strategy will be integrated with per-channel weight quantization. Its resulting Wikitext-2 PPL will be evaluated against baseline schemes that rely on a fixed global $\alpha$ for the entire model.

\subsection{Evaluating Robustness Across the Scale Threshold}
Prior literature establishes that activation outliers become aggressively more prominent and severe as models cross the 6B parameter threshold[]. Therefore the third objective is to evaluate the behavior of the proposed quantile-based and the per-layer migration methods at this critical scale. By measuring Wikitext-2 PPL on OPT-6.7B and OPT-13B, this objective seeks to verify that the proposed methodology scales robustly and withstands the specific regime where baseline methods are most vulnerable to outlier degradation.

\subsection{Verifying Memory Footprint}
The final objective is to verify that the proposed algorithmic refinements preserve the hardware advantages of standard INT8 precision. Because the quantile calculations and per-layer $\alpha$ adjustments are designed as offline modifications to the calibration step, they should incur zero additional compute or memory overhead during deployment. This will be validated by measuring peak VRAM usage and the activation memory footprint during real INT8 inference on GPU hardware, ensuring the deployed models reliably match the theoretical memory reductions of the upstream baseline.

\section{Thesis outline}

The four objectives above are pursued across the remainder of this thesis. Chapter 2 presents the necessary background on Transformer architectures, integer quantization, and the SmoothQuant method that this work builds upon. Chapter 3 describes the proposed method in detail. It walks through per-channel weight quantization, quantile-based smoothing, and the per-layer migration strength as three connected refinements to the baseline. Chapter 4 reports the experimental setup and the empirical results. It covers WikiText-2 perplexity across five OPT sizes, the cross-architecture verification on Falcon-7B and Llama-2-7B, and the real INT8 memory footprint measurement. Chapter 5 concludes the thesis and outlines directions for future work.
